% !TEX program = xelatex
% mycv v4.2.1 mathematical reference
\documentclass[11pt,a4paper]{article}

\usepackage{fontspec}
\setmainfont{Latin Modern Roman}
\setsansfont{Latin Modern Sans}
\setmonofont{Latin Modern Mono}
\usepackage[a4paper,margin=25mm,headheight=15pt]{geometry}
\usepackage{amsmath,amssymb,mathtools,bm}
\usepackage{booktabs,tabularx,array}
\usepackage{xcolor}
\usepackage{microtype}
\usepackage{enumitem}
\usepackage{fancyhdr}
\usepackage{titlesec}
\usepackage{hyperref}

\definecolor{ink}{HTML}{17243A}
\definecolor{blue}{HTML}{2457A6}
\definecolor{muted}{HTML}{536276}
\definecolor{rule}{HTML}{D8E0EA}
\hypersetup{colorlinks=true,linkcolor=blue,urlcolor=blue,citecolor=blue,
  pdftitle={mycv 4.2.1: Mathematical Foundations},
  pdfauthor={Abodunrin Charles Toluwanimi},pdfsubject={Mathematical reference for mycv}}
\titleformat{\section}{\Large\bfseries\color{ink}}{\thesection}{0.65em}{}
\titleformat{\subsection}{\large\bfseries\color{blue}}{\thesubsection}{0.65em}{}
\titleformat{\subsubsection}{\normalsize\bfseries\color{ink}}{\thesubsubsection}{0.65em}{}
\titlespacing*{\section}{0pt}{2.2ex plus .5ex}{1ex}
\titlespacing*{\subsection}{0pt}{1.7ex plus .4ex}{.6ex}
\setlist{nosep,leftmargin=1.5em}
\setlength{\parskip}{0.45em}
\setlength{\parindent}{0pt}
\renewcommand{\arraystretch}{1.2}
\pagestyle{fancy}
\fancyhf{}
\fancyhead[L]{\small\color{muted}mycv 4.2.1 \;|\; Mathematical Foundations}
\fancyhead[R]{\small\color{muted}A. C. Toluwanimi}
\fancyfoot[C]{\small\color{muted}\thepage}
\renewcommand{\headrulewidth}{0.4pt}
\renewcommand{\headrule}{\hbox to\headwidth{\color{rule}\leaders\hrule height \headrulewidth\hfill}}

\DeclareMathOperator{\tr}{tr}
\DeclareMathOperator{\rank}{rank}
\DeclareMathOperator{\softmax}{softmax}
\DeclareMathOperator{\argmax}{arg\,max}
\newcommand{\R}{\mathbb{R}}
\newcommand{\Z}{\mathbb{Z}}
\newcommand{\E}{\mathbb{E}}
\newcommand{\1}{\mathbf{1}}
\newcommand{\I}{\mathbf{I}}
\newcommand{\mycv}{\texttt{mycv}}

\begin{document}
\begin{titlepage}
  \centering
  \vspace*{22mm}
  {\Large\sffamily\bfseries\color{blue} mycv\par}
  \vspace{6mm}
  {\huge\bfseries\color{ink} Mathematical Foundations\par}
  \vspace{4mm}
  {\Large A Pure-NumPy Computer Vision Library\par}
  \vspace{13mm}
  {\large Abodunrin Charles Toluwanimi\par}
  \vspace{2mm}
  {\normalsize Department of Electrical and Electronic Engineering\\
  Obafemi Awolowo University, Ile-Ife\par}
  \vfill
  \begin{minipage}{0.82\textwidth}
  \centering\color{muted}
  A mathematical and implementation reference for the algorithms exposed by
  \mycv{} v4.2.1. This edition extends the earlier v2.0.0 and v3.0.0 notes to
  cover the current library, including tracking, calibration, region geometry,
  neural networks, and the bundled shape classifier.
  \end{minipage}
  \vfill
  {\large Version 4.2.1\par}
  \vspace{2mm}
  {\normalsize September 2026\par}
\end{titlepage}

\tableofcontents
\clearpage

\section{Scope, notation, and implementation conventions}

This monograph follows the current source tree and public API of \mycv{} v4.2.1.
It is a companion to the code, not a claim that every algorithm is a production
vision system. The library implements its core image algorithms with NumPy and
does not require OpenCV, SciPy, or scikit-image. Optional display, image I/O,
and video streaming use separate dependencies.

An image domain is
\[
  \Omega=\{0,\ldots,H-1\}\times\{0,\ldots,W-1\}.
\]
We write a pixel coordinate as $(x,y)$, where $x$ is the column and $y$ is the
row. Array indexing is therefore \texttt{image[y, x]}; rows increase downward. RGB
images have shape $(H,W,3)$ and grayscale images have shape $(H,W)$. Most
photometric and mask functions use 8-bit values, while geometric and statistical
intermediates use floating point.

For a binary mask $M$, foreground means $M(y,x)>0$. Bounding boxes use the
exclusive convention $(y_1,x_1,y_2,x_2)$, so the corresponding crop is exactly
\texttt{image[y1:y2, x1:x2]}. Coordinate order and endpoint convention matter when
passing detections between modules.

\subsection{Array values, axes, and output contracts}

An RGB tensor is stored channel-last: $I\in\{0,\ldots,255\}^{H\times W\times3}$.
This agrees with the way a conventional NumPy view indexes one pixel as
\texttt{I[y,x,:]}. A grayscale frame is $G\in\{0,\ldots,255\}^{H\times W}$.
Several algorithms accept floating-point arrays as well, but each function
documents its expected numeric range and output type. A value range is part of
the mathematical meaning: HSV value is normalized to $[0,1]$, while an 8-bit
mask represents foreground by 255 and background by 0.

The coordinate pair used in geometry is $(x,y)$, but shapes of arrays are
written $(H,W)$. Conversion between the two is not cosmetic: a reported center
$(x,y)$ must be indexed as \texttt{frame[y,x]}. The library's exclusive bounding-box
convention is also deliberate. For a set of foreground coordinates, the
minimal enclosing box is
\[
 y_1=\min_i y_i,\quad y_2=\max_i y_i+1,\quad
 x_1=\min_i x_i,\quad x_2=\max_i x_i+1.
\]
Consequently, \texttt{x2-x1} and \texttt{y2-y1} equal the pixel width and height, and the
NumPy slice includes every foreground pixel without an off-by-one correction.

\subsection{Vectorized work and algorithmic complexity}

Vectorization changes where loops execute; it does not change the mathematical
operator. A broadcast or reduction expresses work over many indices at once,
while NumPy performs the low-level loops. If an image has $N=HW$ pixels, a
pointwise transform is $O(N)$. A direct $k_H\times k_W$ filter is
$O(Nk_Hk_W)$, independent of whether the patch contraction is written with a
Python loop or \texttt{einsum}. Stride tricks create views rather than copies, but
downstream elementwise operations can still allocate storage proportional to
the entire patch tensor.

The implementation uses explicit loops when the dependency structure is
sequential or data-dependent: connected-component equivalence resolution,
contour tracing, RANSAC hypothesis sampling, and greedy NMS. These are part of
the algorithmic design, not accidental failures to vectorize.

\begin{center}
\begin{tabularx}{\textwidth}{@{}l X@{}}
\toprule
Module & Mathematical area and main operations \\
\midrule
\texttt{core}, \texttt{color} & Pointwise maps, luminance, histograms, RGB/HSV coordinates \\
\texttt{filters}, \texttt{morphology} & Discrete operators, set morphology, connected components \\
\texttt{geometry}, \texttt{calibration} & Interpolation, affine/projective maps, homography estimation \\
\texttt{features}, \texttt{detection} & Local moments, shape descriptors, template scores, box overlap \\
\texttt{tracking} & Segmentation, state estimation, data association \\
\texttt{nn}, \texttt{shape\_cnn} & Reverse-mode differentiation and learned shape classification \\
\bottomrule
\end{tabularx}
\end{center}

\section{Point operations and colour representations}

\subsection{RGB to grayscale}

For a pixel with RGB components in $[0,255]$, \mycv{} applies the BT.601-style
luminance map
\[
  Y=0.299R+0.587G+0.114B.
\]
The weights sum to one, keeping the real-valued result within the channel range.
The implementation computes in floating point, clips to $[0,255]$, rounds to
the nearest integer, and returns \texttt{uint8}. The output is a luma-like grayscale
image; it is not a full colourimetric conversion for arbitrary transfer
functions or camera profiles.

The conversion is a linear functional applied independently at every pixel:
\[
 \bm{w}^{\mathsf T}\bm{c}
 =\begin{bmatrix}0.299&0.587&0.114\end{bmatrix}
  \begin{bmatrix}R\\G\\B\end{bmatrix}.
\]
Since the coefficients are nonnegative and sum to one, the unrounded result is
a convex combination of the channels. For white it remains 255; for black it
remains zero. This practical luminance approximation does not linearize
gamma-encoded RGB before weighting.

\subsection{Binary threshold}

For threshold $\tau$, the binary image is
\[
  T_\tau(I)(y,x)=255\,\1\{I(y,x)\geq\tau\}.
\]
The comparison is inclusive. This pointwise segmentation rule is inexpensive,
but a single global threshold can fail under nonuniform illumination or when
foreground and background intensity distributions overlap.

For fixed $\tau$, thresholding is idempotent:
$T_\tau(T_\tau(I))=T_\tau(I)$, because only 0 and 255 remain after one pass.
It is also monotone in the image: if $I_1\le I_2$ pointwise, then
$T_\tau(I_1)\le T_\tau(I_2)$. These are useful algebraic properties of the
mask operator; they do not ensure the mask separates the intended objects.

\subsection{Histogram equalization}

Let $h(k)$ count pixels at intensity $k$ and let $N=HW$. The empirical
probability mass and cumulative distribution are
\[
  p(k)=\frac{h(k)}{N},\qquad F(k)=\sum_{j=0}^{k}p(j).
\]
The lookup table used by the implementation is
\[
  L(k)=\operatorname{round}(255F(k)),\qquad I_{\rm eq}(y,x)=L(I(y,x)).
\]
This monotone remapping can spread occupied intensity levels and improve
contrast. A discrete image does not in general become exactly uniform, and
equalization can also amplify noise in sparsely populated ranges.

The empirical CDF is monotone because each additional bin contributes
nonnegative mass. Therefore the lookup table preserves intensity order up to
quantization ties. For a continuous random variable, the probability integral
transform makes $F(I)$ uniform on $[0,1]$; a finite discrete image only
approximates this ideal. This implementation maps the full CDF and does not
subtract its first nonzero value, so the lowest occupied level need not map to
zero. In particular, a constant image maps to 255 under the implemented rule.

\subsection{RGB to HSV and colour masks}

Normalize $(R,G,B)$ by 255 and set $V=\max(R,G,B)$,
$m=\min(R,G,B)$, and $\Delta=V-m$. Then
\[
 S=\begin{cases}\Delta/V,&V>0,\\0,&V=0,\end{cases}
\]
and hue is chosen from the maximal channel:
\[
H=\begin{cases}
60\bigl((G-B)/\Delta\bmod 6\bigr),&V=R,\\
60\bigl((B-R)/\Delta+2\bigr),&V=G,\\
60\bigl((R-G)/\Delta+4\bigr),&V=B,\\
0,&\Delta=0.
\end{cases}
\]
The returned channels are $H\in[0,360)$, $S,V\in[0,1]$. Hue is undefined
for achromatic colours and is set to zero by convention. A standard colour mask
is an inclusive range test over H, S, and V. \texttt{color\_mask\_hue\_wrap} handles a
hue interval crossing the $0/360$ boundary. \texttt{color\_mask\_adaptive} uses
saturation/value thresholds for near-achromatic samples, where hue is unstable.

The hue branches correspond to sectors of the RGB cube. When red is maximal,
the formula traverses magenta, red, and yellow; the green and blue branches are
offset by 120 and 240 degrees. Modulo arithmetic wraps negative intermediate
angles into $[0,360)$. If $\Delta=0$, all channels are equal and hue has no
geometric meaning. The implementation uses a safe denominator during the
calculation and then sets hue to zero. Boolean masks give deterministic branch
selection when channels tie for the maximum.

An ordinary hue interval $[h_{\min},h_{\max}]$ is tested with a conjunction.
An interval such as [350,10] crosses the circular seam and instead uses
$H\ge350\ \lor\ H\le10$. Sorting the endpoints would select almost the entire
hue circle, so wrapped intervals need their own predicate.

\section{Spatial filters and binary morphology}

\subsection{Two-dimensional convolution}

For image $I$ and kernel $K$, true discrete convolution is
\[
  (I*K)[y,x]=\sum_{i,j}I[y-i,x-j]K[i,j].
\]
Thus the kernel is reflected in both axes before it is applied. \texttt{convolve2d}
uses zero padding for \texttt{same} output and no padding for \texttt{valid} output. A
sliding-window view exposes overlapping image patches and \texttt{einsum} contracts
each patch with the flipped kernel. The view avoids copying the input patches,
although the computation still scales with the number of output positions times
the kernel area.

For a kernel of size $k_H\times k_W$, \texttt{valid} output has shape
$(H-k_H+1,W-k_W+1)$. In \texttt{same} mode, padding by
$p_H=\lfloor k_H/2\rfloor$ and $p_W=\lfloor k_W/2\rfloor$ on both sides gives
the input shape for odd kernels. With an even kernel, this symmetric padding
produces one additional output element along an even-sized dimension. Odd
kernels are therefore the unambiguous choice when exact same-shape centering
matters.

Convolution differs from cross-correlation by reflecting the kernel:
\[
 (I\star K)[y,x]=\sum_{i,j}I[y+i,x+j]K[i,j],\qquad
 (I*K)=I\star\operatorname{flip}(K).
\]
For a symmetric kernel both definitions agree. For a directional derivative
kernel, reflection changes the sign, so flipping is required to implement the
convolution convention. Zero padding assumes the image is zero outside its
domain; near a border this can create an artificial gradient even on a
constant nonzero image.

\subsection{Sobel gradients}

The horizontal Sobel kernel and its transpose are
\[
K_x=\begin{bmatrix}-1&0&1\\-2&0&2\\-1&0&1\end{bmatrix},
\qquad K_y=K_x^{\mathsf T}.
\]
With $G_x=I*K_x$ and $G_y=I*K_y$, gradient magnitude and direction are
\[
  M=\sqrt{G_x^2+G_y^2},\qquad \theta=\operatorname{atan2}(G_y,G_x).
\]
The API returns raw float64 gradient maps and direction, plus a magnitude map
normalized by that image's maximum to the 8-bit range. Consequently, normalized
magnitudes from different images are not directly comparable as absolute
gradient strengths.

The $[1,2,1]$ coefficients smooth across the derivative direction, while
$[-1,0,1]$ estimates a centered first difference. Their outer product gives
$K_x$; transposing exchanges horizontal and vertical roles. A separable
implementation could apply two one-dimensional filters and reduce arithmetic,
but this implementation uses the full 3 by 3 kernels for clarity. Direction
uses \texttt{atan2}, so it distinguishes all quadrants; the zero-gradient case has no
preferred direction and is returned as zero by the numerical convention.

\subsection{Binary morphology as set operations}

Represent a mask by its foreground set $F\subseteq\Omega$, and a structuring
element by $B\subseteq\Z^2$. Dilation and erosion are
\[
 F\oplus B=\{z:\exists b\in B,\ z-b\in F\},\qquad
 F\ominus B=\{z:\forall b\in B,\ z+b\in F\}.
\]
Dilation is an existential (logical ANY) neighborhood reduction; erosion is a
universal (logical ALL) reduction. Opening and closing compose them:
\[
  F\circ B=(F\ominus B)\oplus B,\qquad
  F\bullet B=(F\oplus B)\ominus B.
\]
The code uses constant padding: dilation treats outside pixels as background,
while erosion pads with foreground. That erosion boundary choice avoids
automatically eroding objects merely because they touch the image edge. The
structuring element's active cells alone participate in the reductions.

For a symmetric structuring element, dilation and erosion are dual under
complement: $(F\oplus B)^c=F^c\ominus B$ (with the reflected element in the
general case). On the infinite lattice, opening is anti-extensive and
idempotent, while closing is extensive. The finite image boundary and the
library's padding rules affect these identities at image edges, so the
infinite-lattice laws must be interpreted relative to the padded domain.

With a 3 by 3 all-ones element, dilation can connect foreground pixels at
Chebyshev distance one. Repeating dilation $n$ times expands by the repeated
Minkowski sum $F\oplus nB$. Erosion removes a pixel unless every active
neighborhood cell remains foreground. The kernel is a binary support mask:
nonzero entries are active, and their numeric magnitudes do not weight the
result.

\texttt{grayscale\_dilate} instead computes a sliding maximum and pads with $-\infty$;
it is used to identify local maxima, for example in corner-response maps.

\subsection{Connected components}

Connected-component labelling partitions the foreground into maximal connected
sets under either 4-neighbor or 8-neighbor adjacency. The implementation uses
a raster scan with provisional labels and union-find equivalence resolution,
then a second pass maps roots to consecutive component IDs. Unlike many
pixelwise operations in the package, this algorithm has explicit sequential
loops. For component $C_i$, useful measurements include
\[
 A_i=|C_i|,\quad
 (\bar{x}_i,\bar{y}_i)=\left(\frac{1}{A_i}\sum_{(x,y)\in C_i}x,
                            \frac{1}{A_i}\sum_{(x,y)\in C_i}y\right).
\]
Component properties also expose an exclusive bounding box and optional mean
colour; \texttt{select\_component} can choose the largest component or one near a
requested point.

In the first scan, only already visited causal neighbors can carry a label. If
none is foreground, a new provisional label is created. If several labels
touch the current pixel, their equivalence is recorded in a disjoint-set
forest. Path compression and union by rank keep equivalence operations nearly
constant amortized cost; the scan itself is linear in the number of pixels.
Connectivity changes topology: two diagonal pixels are separate under
4-connectivity but joined under 8-connectivity. Component IDs are identifiers,
not an ordering by area.

\section{Geometry and projective calibration}

\subsection{Bilinear interpolation}

For source coordinate $(x_s,y_s)$, put $x_0=\lfloor x_s\rfloor$,
$x_1=x_0+1$, $\alpha=x_s-x_0$ and similarly $y_0,y_1,\beta$. Then
\begin{align*}
I(x_s,y_s)={}&(1-\alpha)(1-\beta)I[y_0,x_0]
+\alpha(1-\beta)I[y_0,x_1]\\
&+(1-\alpha)\beta I[y_1,x_0]
+\alpha\beta I[y_1,x_1].
\end{align*}
The four weights are nonnegative and sum to one. Coordinates outside the image
are clamped to the nearest edge pixel (replication padding).

The formula is the tensor product of two one-dimensional linear interpolants.
For $(\alpha,\beta)=(0.25,0.5)$, the weights in row-major order are 0.375,
0.125, 0.375, and 0.125. They sum to one, so a constant image remains constant
at fractional coordinates. Interpolation is exact for functions bilinear
inside a cell and approximate for a general sampled image. Clamping the four
integer indices implements edge replication, not constant-zero extension.

\subsection{Rotation and perspective warp}

Using homogeneous column vectors $\tilde p=(x,y,1)^{\mathsf T}$, rotation about
the image center $c=(c_x,c_y)$ is represented by
\[
 M=T(c)R(\theta)T(-c),\qquad
 R(\theta)=\begin{bmatrix}\cos\theta&-\sin\theta&0\\
 \sin\theta&\cos\theta&0\\0&0&1\end{bmatrix}.
\]
To avoid holes, the implementation maps each destination pixel backward through
$M^{-1}$ and samples the source with bilinear interpolation. Rotation is
counter-clockwise in Cartesian coordinates; image rows themselves increase
downward, so visual orientation follows the coordinate convention used by the
implementation.

A homography $H$ maps source points to destination points up to scale:
\[
  \tilde p'\sim H\tilde p.
\]
For each destination pixel, \texttt{warp\_perspective} applies $H^{-1}$ and dehomogenizes
$(u,v,w)^{\mathsf T}$ as $(u/w,v/w)$. It then bilinearly samples the source.
The output canvas is specified by the caller; this function does not infer a
canvas that encloses the warped image.

Backward mapping is a sampling strategy, not a reversal of the requested
geometric transform. If a forward map sends source points to destination
points, then each destination lattice point uses the inverse map to find its
source sample. Every output pixel therefore receives a value, avoiding holes
that can occur when source pixels are pushed forward between destination
pixels. The homogeneous coordinate $w$ must be nonzero. Values near zero are
guarded against division by zero, but samples close to the projective horizon
can remain numerically unstable.

\subsection{Normalized DLT and RANSAC}

For correspondence $(x_i,y_i)\leftrightarrow(x'_i,y'_i)$, the projective
constraint yields two homogeneous linear equations in $h=\operatorname{vec}(H)$:
\[
\begin{bmatrix}-x_i&-y_i&-1&0&0&0&x'_ix_i&x'_iy_i&x'_i\\
0&0&0&-x_i&-y_i&-1&y'_ix_i&y'_iy_i&y'_i\end{bmatrix}h=0.
\]
Stacking $N\ge4$ correspondences gives $Ah=0$. Before building $A$, Hartley
normalization translates each point set to zero centroid and scales it so its
mean radial distance is $\sqrt{2}$. SVD supplies the right singular vector
associated with the smallest singular value. The normalized estimate is
denormalized as
\[
  H=T_{\rm dst}^{-1}H_nT_{\rm src}.
\]
The implementation sets $H_{33}=1$ after denormalization; this assumes that
entry is numerically usable.

For a candidate $H$, the reprojection error is
\[
 e_i=\left\|\pi(H\tilde p_i)-p'_i\right\|_2,
\]
where $\pi(u,v,w)=(u/w,v/w)$. RANSAC repeatedly fits from four sampled pairs,
marks $e_i<\epsilon$ as inliers, retains the largest consensus set, and refits
on all selected inliers. It is robust to outliers when a sufficient inlier
consensus exists; it does not resolve degenerate or poorly distributed point
configurations automatically.

For one correspondence, the cross-product constraint
$\tilde p'_i\times(H\tilde p_i)=0$ supplies three scalar equations, of which
two are independent. Stacking two rows per pair produces a $2N\times9$ matrix.
The homogeneous problem has a nonzero solution because the scale of $H$ is
arbitrary; the unit-norm singular vector fixes scale during the SVD solve.
Hartley normalization reduces coordinate magnitude disparity before products
such as $x'_ix_i$ are formed, improving conditioning for large pixel values.

The minimal sample contains four point pairs, but the final fit should use all
inliers. Collinear or nearly collinear points do not constrain a general
projective map well. A small inlier threshold rejects legitimate noisy pairs;
a large one admits outliers. In an independent-sampling model, if each pair is
an inlier with probability $w$, a four-pair sample is all-inlier with
probability $w^4$. To obtain success probability $p$, a conventional bound is
\[
 N_{\rm iter}\ge\frac{\log(1-p)}{\log(1-w^4)}.
\]
The implementation uses the requested iteration count rather than adapting
this bound during sampling.

\section{Features and object measurements}

\subsection{Harris corners}

For image gradients $I_x,I_y$, form local Gaussian-weighted second moments
\[
 S_{xx}=G_\sigma*I_x^2,\quad S_{xy}=G_\sigma*(I_xI_y),\quad
 S_{yy}=G_\sigma*I_y^2,
 \qquad
 M=\begin{bmatrix}S_{xx}&S_{xy}\\S_{xy}&S_{yy}\end{bmatrix}.
\]
The Harris response is
\[
  R=\det(M)-k\,\tr(M)^2.
\]
Large positive responses indicate variation in two directions and are corner
candidates; edge-like regions have one dominant eigenvalue. \texttt{detect\_harris\_corners}
adds thresholding, local-maximum filtering, and optional mask restriction.

The matrix is positive semidefinite because it is a weighted sum of gradient
outer products. Let its eigenvalues be $\lambda_1,\lambda_2\ge0$. The response
can be written
\[
 R=\lambda_1\lambda_2-k(\lambda_1+\lambda_2)^2.
\]
If one eigenvalue is small, the neighborhood changes mostly along one
direction and is edge-like. If both are large, it changes in two independent
directions and is corner-like. If both are small, it is locally flat. The
parameter $k$ balances determinant against total gradient energy; it is not an
absolute threshold because response scale also depends on image contrast and
smoothing.

Non-maximum suppression compares each response with a local maximum map. A
threshold removes weak responses, while the neighborhood size determines how
closely spaced two corners can be while both survive. Equal-valued plateaus
can retain several neighboring maxima, so the returned points are candidates
rather than guaranteed unique physical corners.

\subsection{Hough lines}

A line is represented in normal form
\[
  \rho=x\cos\theta+y\sin\theta.
\]
For each edge pixel and sampled $\theta$, the implementation computes a $\rho$
bin and accumulates a vote. Peaks in the accumulator are line hypotheses;
\texttt{count\_hough\_lines} clusters nearby peaks to estimate a count. The result
depends on angle/rho resolution, edge quality, and peak thresholds.

The normal form uses a unit normal vector
$n(\theta)=(\cos\theta,\sin\theta)$, so $\rho=n^{\mathsf T}(x,y)$ is signed
distance from the origin. A single image point traces a sinusoid in the
$(\theta,\rho)$ accumulator. Points lying on the same image line produce
sinusoids that intersect at the shared line parameters. Restricting $\theta$
to an interval of length $\pi$ avoids representing an unoriented line twice,
because $(\rho,\theta)$ and $(-\rho,\theta+\pi)$ are equivalent.

The finite accumulator discretizes both parameters. Coarse bins merge nearby
lines; fine bins spread votes and make peaks sensitive to pixel noise. The
maximum absolute $\rho$ follows from the image diagonal. Indexed addition is
needed because many edge pixels may vote for the same bin. One physical line
often creates neighboring high-vote bins, which is why line counting clusters
peaks instead of counting every thresholded cell.

\subsection{Region moments and shape descriptors}

For foreground coordinates $(x_i,y_i)$, area and centroid are
\[
 A=\sum_i1,\quad \bar{x}=A^{-1}\sum_i x_i,\quad
 \bar{y}=A^{-1}\sum_i y_i.
\]
The centered second moments $\mu_{20}=\sum_i(x_i-\bar{x})^2$,
$\mu_{02}=\sum_i(y_i-\bar{y})^2$, and
$\mu_{11}=\sum_i(x_i-\bar{x})(y_i-\bar{y})$ define the covariance-like matrix
\[
 C=A^{-1}\begin{bmatrix}\mu_{20}&\mu_{11}\\\mu_{11}&\mu_{02}\end{bmatrix}.
\]
Its principal orientation is
\[
 \phi=\tfrac12\operatorname{atan2}(2\mu_{11},\mu_{20}-\mu_{02}),
\]
and with eigenvalues $\lambda_1\ge\lambda_2\ge0$, eccentricity is
$e=\sqrt{1-\lambda_2/\lambda_1}$ (defined as zero for degenerate spread).

The exclusive bounding box has width $w=x_2-x_1$ and height $h=y_2-y_1$.
The library distinguishes true foreground area from rectangle area:
\[
 \text{extent}=A/(wh),\quad \text{aspect ratio}=w/h.
\]
With traced perimeter $P$ and convex hull area $A_h$, it also reports
\[
 \text{circularity}=\min\!\left(1,\frac{4\pi A}{P^2}\right),\qquad
 \text{solidity}=A/A_h.
\]
These are discrete estimators and can be sensitive to segmentation and small
objects. The heuristic classifier calls a region a circle only when both
circularity is at least 0.85 and extent is at most 0.88; otherwise it uses
aspect ratio to separate square-like from rectangular shapes. Colour is chosen
as the nearest of six fixed RGB reference colours.

The covariance matrix is the second central-moment matrix divided by the
foreground count. Its principal eigenvector gives the major axis, and
$\tfrac12\operatorname{atan2}(2\mu_{11},\mu_{20}-\mu_{02})$ is its closed-form
orientation. Eccentricity is zero for equal eigenvalues and approaches one as
the minor-axis spread vanishes. Orientation is ambiguous for nearly isotropic
regions; the implementation explicitly returns zero when anisotropic terms
are numerically negligible.

The perimeter estimator traces an 8-connected boundary with a Moore-neighbor
walk and sums chain-code step lengths: one for orthogonal steps and $\sqrt2$
for diagonal steps. Circularity uses this perimeter and foreground pixel
count, then clips to one to contain rasterization overshoot. Solidity compares
pixel count with the polygon area of the convex hull. A degenerate hull can
have zero area, in which case solidity is unavailable.

The heuristic requires two signals for a circle because circularity alone can
misclassify small rasterized squares. Extent tends to be one for a filled
axis-aligned rectangle, while a circle occupies only part of its enclosing
box. Oblique rectangles, holes, and imperfect segmentation change these
descriptors, so the classification is a rule-based estimate rather than a
shape guarantee.

\section{Template detection and box suppression}

\subsection{Normalized cross-correlation}

At each valid image patch $P_{xy}$, let $\tilde P=P_{xy}-\bar P$ and
$\tilde T=T-\bar T$. The normalized cross-correlation score is
\[
 \operatorname{NCC}(x,y)=
 \frac{\frac{1}{n}\sum_j\tilde P_j\tilde T_j}
 {\sigma_P\sigma_T},\qquad n=k_Hk_W.
\]
This is Pearson correlation over patch entries and lies in $[-1,1]$. It is
invariant to additive offsets and positive contrast scaling in the ideal
non-degenerate case. If the patch or template has near-zero variance, the
implementation returns zero. The output map has shape
$(H-k_H+1,W-k_W+1)$, corresponding to each patch's top-left location. A memory
guard prevents materializing excessively large patch-statistic temporaries.

Flatten each centered patch and template into vectors $p,t\in\R^n$. Since
$\|p\|_2=\sqrt{n}\sigma_P$ and $\|t\|_2=\sqrt{n}\sigma_T$, the factors of
$n$ cancel and
\[
 \operatorname{NCC}=\frac{p^{\mathsf T}t}{\|p\|_2\|t\|_2}.
\]
Thus NCC is cosine similarity after centering. If an image patch is transformed
to $aP+b$ with $a>0$, its centered vector becomes $ap$; normalization cancels
$a$, and the offset $b$ disappears. A negative $a$ reverses the score sign.
This explains the score's brightness and contrast behavior without implying
invariance to spatial deformation, rotation, occlusion, or a nonlinear tone
curve.

The numerator uses the mean of elementwise products and the denominator uses
population standard deviations over the patch cells. Flat patches and flat
templates carry no spatial pattern, so their undefined zero-over-zero scores
are set to zero. The result is clipped to $[-1,1]$ against floating-point
roundoff. If there are $M=(H-k_H+1)(W-k_W+1)$ candidate positions and a template
has $q=k_Hk_W$ cells, the direct patch calculation costs $O(Mq)$ arithmetic.
The patch view itself is cheap, but centered-patch arrays and products can
materialize $O(Mq)$ storage. A memory guard fails early rather than attempting
an allocation likely to exhaust memory.

\subsection{Pyramids and multi-scale matching}

Before downsampling by two, a 5 by 5 binomial kernel is applied:
\[
 G=\frac{1}{256}\begin{bmatrix}1&4&6&4&1\end{bmatrix}^{\mathsf T}
 \begin{bmatrix}1&4&6&4&1\end{bmatrix},\qquad
 I_{\ell+1}=(I_\ell*G)[::2,::2].
\]
The low-pass step attenuates frequencies that would alias under decimation.
Multi-scale template matching searches pyramid levels, mapping detections back
to the original image coordinates. It trades fine localization and small-object
coverage for less work at coarse levels.

The binomial filter is separable and sums to one, so it preserves a constant
image away from the zero-padded border. Downsampling every second row and
column halves each spatial dimension approximately; at level $\ell$ the shape
is $\lceil H/2^\ell\rceil\times\lceil W/2^\ell\rceil$. Across an unlimited
pyramid, total pixel storage is bounded by
\[
 HW\left(1+\frac14+\frac1{16}+\cdots\right)=\frac43HW.
\]
Blur-before-decimation suppresses high frequencies above the new Nyquist
limit. It reduces aliasing but does not perfectly eliminate it because the
finite binomial kernel is only an approximate low-pass filter.

At level $\ell$, a one-pixel displacement corresponds to roughly $2^\ell$
pixels at the original scale. Coarse-level matching can therefore locate a
large object cheaply, followed by finer search for localization. Very small
objects can disappear at coarse levels, and pyramid quantization shifts mapped
coordinates by up to a small fraction of a level's scale.

\subsection{Non-maximum suppression}

For exclusive-coordinate boxes $A,B$, intersection area is
\[
 |A\cap B|=\max(0,y_2^I-y_1^I)\max(0,x_2^I-x_1^I),
\]
where $y_1^I=\max(y_1^A,y_1^B)$, $y_2^I=\min(y_2^A,y_2^B)$, and likewise for
$x$. Intersection over union is
\[
 \operatorname{IoU}(A,B)=\frac{|A\cap B|}{|A|+|B|-|A\cap B|}.
\]
Greedy NMS repeatedly keeps the highest-scoring remaining box and removes
boxes whose IoU exceeds the selected threshold. Lower thresholds suppress more
overlapping candidates.

For candidate boxes sorted by descending score, greedy NMS keeps the first,
removes candidates with $\operatorname{IoU}>\tau$, and repeats on the survivors.
The comparison is strict: a candidate exactly at the threshold is retained.
Suppression is relative to boxes already kept, not a symmetric clustering of
all pairwise overlaps. In the worst case, $N$ boxes require $O(N^2)$ overlap
comparisons, although each current-box comparison is vectorized across the
remaining candidates. A threshold near zero is very aggressive; a threshold
near one retains highly overlapping boxes.

\section{Segmentation and motion tracking}

\subsection{Motion and colour masks}

For consecutive grayscale frames, motion is detected using the absolute
difference
\[
 D_t(y,x)=|I_t(y,x)-I_{t-1}(y,x)|,\qquad
 M_t(y,x)=255\,\1\{D_t(y,x)\ge\tau\}.
\]
The implementation widens unsigned inputs before subtraction to avoid uint8
wraparound. Colour tracking instead thresholds HSV values; morphology can
clean small gaps or isolated foreground pixels before component extraction.

For a binary mask $M$, the foreground centroid returned by
\texttt{calculate\_centroid} is
\[
  (\bar{x},\bar{y})=\left(\frac{\sum xM}{\sum M},
                         \frac{\sum yM}{\sum M}\right),
\]
with the mask interpreted as foreground membership. It returns \texttt{None} for an
empty mask. Disconnected regions should normally be split into components
before tracking so their combined centroid does not fall between objects.

If the binary membership values are $m_i\in\{0,1\}$, the centroid is the
first spatial moment divided by the zeroth moment:
\[
 m_{00}=\sum_i m_i,\quad m_{10}=\sum_i x_i m_i,\quad
 m_{01}=\sum_i y_i m_i,\quad
 (\bar{x},\bar{y})=(m_{10}/m_{00},m_{01}/m_{00}).
\]
Using mask intensities of 255 instead of 1 gives the same result because the
constant factor cancels. A union of two objects produces a weighted average
of their centers, which may lie in background; connected-component splitting
prevents this failure for separated blobs.

Frame differencing responds to object motion, camera motion, lighting flicker,
sensor noise, and compression artifacts. Raising the difference threshold
suppresses small intensity changes but can miss slow or low-contrast objects.
The difference between two silhouettes often highlights changed boundaries
rather than the object's entire interior. Opening and closing can reduce
isolated noise or small mask gaps, but their structuring-element size must
match the image scale.

\subsection{Kalman centroid filter}

The single-object state is $s_t=[x_t,y_t,v_{x,t},v_{y,t}]^{\mathsf T}$. With
time step $\Delta t$, the constant-velocity model is
\[
 s_t=F_ts_{t-1}+w_t,\quad
 F_t=\begin{bmatrix}1&0&\Delta t&0\\0&1&0&\Delta t\\
 0&0&1&0\\0&0&0&1\end{bmatrix},\quad
 z_t=Hs_t+v_t,
 \quad H=\begin{bmatrix}1&0&0&0\\0&1&0&0\end{bmatrix}.
\]
For covariance $P$, process noise $Q$, and measurement noise $R$, prediction is
\[
 \hat{s}^-_t=F_t\hat{s}_{t-1},\qquad P^-_t=F_tP_{t-1}F_t^{\mathsf T}+Q.
\]
With residual $r_t=z_t-H\hat{s}^-_t$ and innovation covariance
$S_t=HP^-_tH^{\mathsf T}+R$, the gain and correction are
\[
 K_t=P^-_tH^{\mathsf T}S_t^{-1},\quad
 \hat{s}_t=\hat{s}^-_t+K_tr_t.
\]
The implementation solves linear systems rather than explicitly inverting the
innovation covariance and uses the Joseph covariance update
\[
 P_t=(I-K_tH)P^-_t(I-K_tH)^{\mathsf T}+K_tRK_t^{\mathsf T},
\]
which better preserves symmetry and positive semidefiniteness under finite
precision. Optional Mahalanobis gating rejects measurements with
$r_t^{\mathsf T}S_t^{-1}r_t$ above a configured threshold. Measurement
confidence can adjust effective measurement noise. The multi-object tracker
associates detections with tracks using spatial distance and manages IDs,
missed frames, and track creation/deletion; it is a lightweight nearest
association scheme rather than a globally optimal assignment solver.

The update can be derived by minimizing posterior uncertainty under a linear
measurement model. The innovation is $r=z-H\hat{s}^-$ and its covariance is
$S=HP^-H^{\mathsf T}+R$. The Kalman gain weights the innovation against prior
uncertainty: if measurement noise grows, the gain shrinks and the estimate
trusts prediction more. This implementation uses diagonal covariances
$Q=qI_4$ and $R=rI_2$; the scalar process-noise setting is not a full
physically derived white-noise acceleration model.

The Joseph form follows by writing posterior error as
$e=(I-KH)e^- - Kv$ and assuming independent prior and measurement noise. Its
covariance expands to
$(I-KH)P^-(I-KH)^{\mathsf T}+KRK^{\mathsf T}$. Unlike the algebraically
equivalent short form $(I-KH)P^-$, this expression is less likely to lose
symmetry or positive semidefiniteness through floating-point rounding. Linear
solves also avoid explicitly forming $S^{-1}$.

Squared Mahalanobis distance is
\[
 d^2=r^{\mathsf T}S^{-1}r.
\]
It measures residual size in units of predicted uncertainty. Under a calibrated
two-dimensional Gaussian innovation, a threshold of 9.21 is approximately the
99th percentile of a chi-squared distribution with two degrees of freedom.
The confidence option scales measurement noise as
$R_{\rm eff}=R_{\rm base}/c$; smaller $c$ makes the filter trust the
observation less.

The filter commits each prediction to state and covariance. During occlusion,
call \texttt{predict()} once per elapsed frame so velocity propagates through time.
On a frame with a measurement, call \texttt{update()} rather than both methods for
that same frame. The first measurement initializes position and sets velocity
to zero; later measurements estimate velocity through recursive correction.

For multiple objects, greedy association computes Euclidean distances between
predicted track centers and new component centroids, sorts candidate pairs by
distance, and accepts a pair if neither endpoint was previously assigned and
the distance is within the configured maximum. Unmatched detections start new
IDs; unmatched tracks advance without a measurement and increment a
missed-frame counter. The pair sorting cost is roughly $O(TD\log(TD))$ for
$T$ tracks and $D$ detections. A locally short assignment may prevent the
globally best pairing; a Hungarian solver would minimize total cost but is not
part of this implementation.

\section{Neural-network toolkit and bundled shape model}

\subsection{Layers and reverse-mode differentiation}

The \texttt{mycv.nn} module supplies convolution, pooling, dense, and ReLU layers,
composed in a sequential model. For a dense layer
\[
 Z=XW+b,\quad \delta=\frac{\partial L}{\partial Z},
\]
the gradients are
\[
 \frac{\partial L}{\partial W}=X^{\mathsf T}\delta,\qquad
 \frac{\partial L}{\partial b}=\sum_{\rm batch}\delta,\qquad
 \frac{\partial L}{\partial X}=\delta W^{\mathsf T}.
\]
For ReLU, $f(a)=\max(0,a)$ and the backward multiplier is $\1\{a>0\}$.
Convolution and pooling implement the corresponding local reverse operations.
The softmax cross-entropy loss for logits $z$ and one-hot target $y$ is
\[
 L(z,y)=-\sum_c y_c\log\softmax(z)_c,
 \qquad \frac{\partial L}{\partial z}=\softmax(z)-y.
\]
The SGD optimizer updates parameters as $\theta\leftarrow\theta-\eta\nabla_\theta L$.

\subsubsection{Convolution, pooling, and tensor layout}

The neural-network layer uses NCHW tensors
$X\in\R^{N\times C_{\rm in}\times H\times W}$. CNN convention calls the
spatial operation cross-correlation: the filter is not flipped. For output
channel $o$ and location $(y,x)$,
\[
 Y_{n,o,y,x}=b_o+\sum_{c,i,j}
 X_{n,c,ys+i,xs+j}W_{o,c,i,j},
\]
where $s$ is stride and the input is padded according to \texttt{same} or \texttt{valid}.
This is intentionally different from the mathematical-convolution convention
in \texttt{filters.convolve2d}.

For a local derivative $\delta=\partial L/\partial Y$, the weight gradient
sums input patches times $\delta$ over batch and output positions, while the
input gradient scatters each $\delta$ back through the corresponding weights.
Bias gradients sum $\delta$ over batch and spatial positions. The code loops
over kernel offsets and uses array contractions across batch, channels, and
spatial positions.

Non-overlapping max pooling retains the largest value in each $k\times k$
window. Its backward pass routes gradient to the winning input; exact ties
divide the gradient evenly across all tied maxima. Average pooling returns the
window mean and divides the upstream gradient evenly among the $k^2$ inputs.
Both pooling layers require height and width divisible by the pool size.
Flatten changes layout from $(N,C,H,W)$ to $(N,CHW)$ and reverses that reshape
during the backward pass.

\subsubsection{Numerical gradient checking}

An analytic derivative can be checked against a centered finite difference.
For parameter $\theta_j$ and small $\epsilon$,
\[
 g_j^{\rm numeric}=\frac{L(\theta_j+\epsilon)-L(\theta_j-\epsilon)}{2\epsilon}.
\]
The relative discrepancy between analytic and numeric gradients should be
small away from nondifferentiable points such as ReLU zero crossings and
max-pool ties. The project tests use finite-difference checks to verify layer
backward passes, which is especially valuable for convolution input-gradient
scattering and pooling tie behavior.

The combined softmax cross-entropy calculation subtracts each row's maximum
logit before exponentiation (the log-sum-exp stabilization). The mean loss over
a batch of $N$ examples is returned, so its gradient includes a factor $1/N$.
The optimizer applies plain stochastic gradient descent; it does not include
momentum, adaptive per-parameter rates, weight decay, or automatic batching.

\subsection{Shape CNN}

\texttt{predict\_shape} uses bundled weights to classify a cropped, non-empty binary
mask as Circle, Square, or Rectangle. The model input is normalized and
resampled to its fixed small input size before inference. Its output contains
the winning label, its softmax confidence, and per-class probabilities. Blank
masks raise \texttt{ValueError}. The model is a narrow shape classifier and augments
the geometric heuristic; it is not a general object detector. Version 4.2.1
retrained the weights with the full square rotation range and applied
degradation augmentation before resizing. The repository metadata records the
training configuration and the measured synthetic stress-test score; that
score should not be interpreted as a guarantee on arbitrary real camera data.

The architecture is fixed to match the bundled parameter file:
\[
 (1\!\to\!8,3\times3)\to\mathrm{ReLU}\to\mathrm{MaxPool}_2\to
 (8\!\to\!16,3\times3)\to\mathrm{ReLU}\to\mathrm{MaxPool}_2
 \to\mathrm{Flatten}\to 576\!\to\!32\to\mathrm{ReLU}\to32\!\to\!3.
\]
Both convolutions use same padding. A $24\times24$ single-channel mask becomes
$12\times12$ and then $6\times6$ after pooling; flattening therefore gives
$16\cdot6\cdot6=576$ inputs to the first dense layer.

Preprocessing thresholds the supplied mask at zero, finds its tight
foreground bounding box, preserves aspect ratio, scales the longer side to 24
pixels, centers the result in a zero canvas, bilinearly samples, then thresholds
at 0.5. This preserves the elongation cue that distinguishes a rectangle from
a square. Empty masks are rejected before inference. The default model loads
bundled weights once and caches them for repeated calls.

The current metadata records 2,500 synthetic training samples per class,
400 validation samples per class, 99.17 percent clean synthetic validation
accuracy, and 96 percent on a combined rotation-and-degradation stress test.
Uniform 3 percent salt-and-pepper flips remain a weak case: 85 percent overall
and 65 percent for circles in the recorded benchmark. These measurements are
distribution-specific; collecting and evaluating real camera masks is needed
before treating confidence values as calibrated for a deployment domain.

\section{Putting the algorithms together}

The modules are designed as composable operators. A useful shape-tracking path
starts with an RGB frame, creates a colour or intensity mask, cleans the mask,
splits it into connected components, measures each component, optionally
classifies a cropped component mask, and passes the resulting centroids to a
multi-object tracker. In symbols, for frame $I_t$ the path can be written
\[
 I_t\xrightarrow{\text{colour threshold}}M_t
 \xrightarrow{\text{morphology}}\widetilde M_t
 \xrightarrow{\text{components}}\{C_{t,i}\}
 \xrightarrow{\text{moments}}\{z_{t,i}\}
 \xrightarrow{\text{association/filter}}\{\hat s_{t,j}\}.
\]
Each arrow changes the kind of information available. A mask records
pixel membership; components separate disconnected regions; region moments
turn pixels into measurements; the Kalman state adds temporal position and
velocity.

\subsection{A single-frame measurement sequence}

For each frame, colour segmentation produces a binary mask. Morphological
opening may remove isolated bright pixels and closing may bridge small gaps.
Connected components then partition the foreground into blobs. A component
area threshold $A_i\ge A_{\min}$ rejects very small candidates. The centroid
is the component's position measurement, while the exclusive bounding box
supplies the correct crop for an object-mask classifier.

\[
 \text{RGB frame}\to\text{HSV mask}\to\text{opened/closed mask}\to
 \{(A_i,\operatorname{bbox}_i,z_i)\}\to\text{shape label and track update}.
\]

The ordering matters. If measurements are taken before components are split,
two separated objects share one centroid. If shape metrics are computed from a
fixed template window rather than the segmented object's mask, extent and
aspect ratio describe the window, not the physical shape. The segmentation
should therefore be restricted to a region of interest around a locator and
then measured on the resulting foreground pixels.

\subsection{A two-frame motion sequence}

Given grayscale frames $G_{t-1}$ and $G_t$, absolute frame differencing
produces $M_t$. After optional morphology and component filtering, each
surviving centroid is a noisy position measurement. A single-object filter
predicts its state, compares the measured centroid against the predicted
position, and updates position and velocity. If no component is detected,
prediction-only evolution carries the state through a short occlusion.

The measurement noise should reflect centroid reliability: a large, clean
blob generally gives a more stable centroid than a tiny noisy component. The
confidence parameter expresses this relationship by scaling $R$. Motion masks
often represent only changed borders, however, so their centroid may not
coincide with the center of the moving object. Colour masks can provide a
better full silhouette when the target colour is separable from the scene.

\subsection{Choosing a locator and a measurer}

Template matching is an appearance-based locator; a colour mask is a
colour-based locator; connected components are region candidates. None alone
guarantees an accurate shape label. \texttt{select\_component} can choose a component
near a template peak, after which region metrics describe the selected mask.
This division keeps the role of each primitive clear:
\begin{itemize}
  \item \textbf{Locate:} identify a region of interest with NCC or colour.
  \item \textbf{Segment:} isolate foreground pixels in or near that region.
  \item \textbf{Measure:} compute area, centroid, orientation, and descriptors.
  \item \textbf{Track:} associate measurements over time and filter their noise.
\end{itemize}

\section{Cost and parameter interpretation}

The following costs describe dominant arithmetic for direct implementations;
they omit constants, allocation overhead, and image-dependent early exits.
\begin{center}
\begin{tabularx}{\textwidth}{@{}l l X@{}}
\toprule
Operation & Dominant cost & Parameters that change the result \\
\midrule
Pointwise transforms & $O(HW)$ & threshold, intensity encoding \\
Direct convolution & $O(HWk_Hk_W)$ & kernel, padding convention \\
Binary morphology & $O(HWk_Hk_W)$ & structuring element and its active support \\
Connected components & $O(HW\,\alpha(HW))$ amortized & 4- or 8-connectivity \\
NCC template map & $O(HWk_Hk_W)$ & template, score threshold, search region \\
Hough voting & $O(E N_\theta)$ & edge pixels, angle and rho resolution \\
RANSAC homography & $O(RN)$ after each four-point fit & iterations, reprojection threshold \\
Greedy box NMS & $O(B^2)$ worst case & IoU threshold and candidate count \\
Kalman update & constant-size matrix algebra & $Q$, $R$, $\Delta t$, gate \\
\bottomrule
\end{tabularx}
\end{center}
Here $E$ is the number of edge pixels, $N_\theta$ is the number of sampled
angles, $R$ is the RANSAC iteration count, $N$ is the number of correspondences,
and $B$ is the number of proposed boxes. The connected-component inverse
Ackermann factor $\alpha$ grows so slowly that the scan is effectively linear.

Parameters have different units. A difference threshold is in grayscale
intensity units; a reprojection threshold is in pixels; an IoU threshold is
dimensionless; and a Mahalanobis gate is a squared standardized distance.
Transferring a numeric value between these settings has no meaning. Thresholds
should be selected from validation examples representative of the deployment
images and object scales.

\section{Numerical practice, limits, and reproducibility}

\begin{itemize}
  \item Cast unsigned images before subtraction or multiplication when signed
        differences or products are required; uint8 arithmetic can wrap.
  \item Use float64 for moments, correlations, projective transforms, and
        covariance updates where dynamic range and precision matter.
  \item Guard divisions at achromatic HSV pixels, flat NCC patches, empty
        masks, and projective points near infinity.
  \item Stride-based patch views are memory-efficient views, but reductions and
        elementwise products may allocate large temporaries. Template matching
        therefore enforces a temporary-memory ceiling by default.
  \item Connected-component labelling, contour tracing, RANSAC iterations,
        and greedy suppression contain sequential algorithmic steps; the
        library is not loop-free even though many array operations are vectorized.
  \item Results depend on image encoding, thresholds, scale, mask quality,
        noise, and camera geometry. The functions expose algorithmic building
        blocks and do not automatically solve these calibration choices.
\end{itemize}

The neural-network layer backward passes are covered by finite-difference
gradient checks in the project tests. Unit tests also exercise the numerical
and shape conventions of the image, tracking, detection, and calibration APIs.
For exact parameter meanings and return-value details, the Python docstrings
remain authoritative for the installed version.

\section{API map}

\begin{center}
\small
\begin{tabularx}{\textwidth}{@{}l X@{}}
\toprule
Module & Public concepts \\
\midrule
\texttt{core} & \texttt{rgb\_to\_grayscale}, \texttt{threshold} \\
\texttt{color} & \texttt{histogram\_equalize}, \texttt{rgb\_to\_hsv} \\
\texttt{filters} & \texttt{convolve2d}, Sobel kernels and edge response \\
\texttt{morphology} & Binary/grayscale morphology, connected components, selection \\
\texttt{geometry} & Bilinear sampling, rotation, perspective warp \\
\texttt{features} & Harris and Hough features, region metrics, shape heuristic \\
\texttt{detection} & NCC matching, pyramids, multiscale matching, NMS \\
\texttt{tracking} & Motion/colour masks, centroid and Kalman trackers \\
\texttt{calibration} & Normalized DLT, reprojection error, RANSAC \\
\texttt{nn}, \texttt{shape\_cnn} & NumPy layers, optimizer, loss, trained shape predictor \\
\texttt{streaming} & Optional PyAV \texttt{StreamReader} \\
\bottomrule
\end{tabularx}
\end{center}

\subsection{Function-by-function guide}

This index describes the purpose of each exported operation. It is deliberately
shorter than the function docstrings, which define exact argument defaults and
validation behavior.

\subsubsection{Core, colour, and filters}
\begin{itemize}
 \item \texttt{rgb\_to\_grayscale(image)} maps an RGB tensor to 8-bit grayscale.
 \item \texttt{threshold(image, tau)} returns a 0/255 mask using the inclusive comparison $I\ge\tau$.
 \item \texttt{histogram\_equalize(image)} applies a cumulative-histogram lookup table to one grayscale image.
 \item \texttt{rgb\_to\_hsv(image)} converts RGB bytes to float32 HSV with hue in degrees.
 \item \texttt{convolve2d(image, kernel, padding)} computes true single-channel convolution with \texttt{same} or \texttt{valid} padding.
 \item \texttt{sobel\_edge\_detection(image)} returns $G_x$, $G_y$, normalized magnitude, and gradient direction.
 \item \texttt{SOBEL\_X}, \texttt{SOBEL\_Y} expose the two 3 by 3 derivative masks.
\end{itemize}

\subsubsection{Morphology and geometry}
\begin{itemize}
 \item \texttt{dilate(image, kernel)} and \texttt{erode(image, kernel)} apply binary ANY and ALL neighborhood reductions.
 \item \texttt{opening(image, kernel)} and \texttt{closing(image, kernel)} compose erosion/dilation in opposite orders.
 \item \texttt{grayscale\_dilate(image, size)} computes a sliding maximum for local-peak filtering.
 \item \texttt{label\_connected\_components(mask, connectivity)} labels separate foreground regions and returns the component count.
 \item \texttt{component\_properties(labels, ...)} measures area, box, centroid, and available per-component statistics.
 \item \texttt{select\_component(...)} selects a component by area or proximity to a requested coordinate.
 \item \texttt{bilinear\_interpolate(image, xs, ys)} samples one grayscale plane at fractional coordinates.
 \item \texttt{rotate\_image(image, angle\_deg)} rotates a grayscale image around its center using backward mapping.
 \item \texttt{warp\_perspective(image, H, output\_shape)} applies a forward homography by inverse destination sampling.
\end{itemize}

\subsubsection{Features and detection}
\begin{itemize}
 \item \texttt{harris\_corner\_response(Gx, Gy, ...)} computes the scalar structure-tensor response map.
 \item \texttt{detect\_harris\_corners(...)} thresholds and locally suppresses that response, optionally within a mask.
 \item \texttt{hough\_line\_transform(edges, ...)} builds the rho/theta vote accumulator and returns candidate peaks.
 \item \texttt{count\_hough\_lines(edges, ...)} clusters accumulator peaks into a line count and related details.
 \item \texttt{convex\_hull(points)} returns the monotone-chain hull of 2D coordinates.
 \item \texttt{extract\_object\_metrics(mask, ...)} computes exclusive box, pixel area, moments, perimeter, circularity, solidity, and optional edge features.
 \item \texttt{draw\_bounding\_box(image, bbox, ...)} draws an exclusive-coordinate box in-place on an RGB array.
 \item \texttt{classify\_object(image, metrics, mask)} selects the nearest reference colour and applies the geometric shape rules.
 \item \texttt{match\_template\_ncc(image, template)} returns a valid-position normalized-correlation response map.
 \item \texttt{find\_template\_matches(map, template\_shape, ...)} thresholds scores, forms exclusive boxes, and applies NMS.
 \item \texttt{gaussian\_pyramid(image, levels)} creates blur-then-decimate image scales.
 \item \texttt{match\_template\_ncc\_multiscale(...)} searches at pyramid scales and maps candidate detections to the base image.
 \item \texttt{non\_max\_suppression(boxes, scores, iou\_threshold)} removes redundant high-overlap proposals.
\end{itemize}

\subsubsection{Tracking and calibration}
\begin{itemize}
 \item \texttt{compute\_motion\_mask(current, previous, threshold)} differences two grayscale frames safely in signed arithmetic.
 \item \texttt{color\_mask(hsv, lower, upper)} selects pixels within a standard HSV interval.
 \item \texttt{color\_mask\_hue\_wrap(...)} accepts a hue interval that crosses 0 degrees.
 \item \texttt{color\_mask\_adaptive(...)} switches near-achromatic samples to saturation/value segmentation.
 \item \texttt{calculate\_centroid(mask)} returns the foreground center or \texttt{None} for an empty mask.
 \item \texttt{TemporalSmoother(window)} provides a ring-buffer temporal average for array frames.
 \item \texttt{kalman\_filter\_predict(state, P, F, Q)} computes one linear state/covariance prediction.
 \item \texttt{kalman\_filter\_update(state, P, z, H, R)} corrects the prediction with a position observation.
 \item \texttt{mahalanobis\_gate(...)} evaluates squared innovation distance against its covariance.
 \item \texttt{KalmanCentroidTracker} stores one position/velocity state, supports per-frame $\Delta t$, missing measurements, gating, and confidence-scaled noise.
 \item \texttt{MultiObjectKalmanTracker} manages track IDs, greedy centroid association, prediction-only misses, and stale-track removal.
 \item \texttt{normalize\_points(points)} returns Hartley-normalized coordinates and the transform matrix.
 \item \texttt{solve\_homography\_dlt(src, dst)} estimates a normalized-DLT homography from at least four pairs.
 \item \texttt{reprojection\_error(H, src, dst)} returns one Euclidean pixel error per correspondence.
 \item \texttt{ransac\_homography(src, dst, ...)} returns an inlier-refined matrix and final Boolean inlier mask.
\end{itemize}

\subsubsection{Neural networks, shape classification, and streaming}
\begin{itemize}
 \item \texttt{mycv.nn.Conv2D} computes a trainable NCHW cross-correlation layer; \texttt{MaxPool2D} and \texttt{AvgPool2D} downsample non-overlapping windows.
 \item \texttt{Flatten}, \texttt{Dense}, and \texttt{ReLU} provide reshape, affine, and elementwise activation layers.
 \item \texttt{Sequential} composes layers and exposes forward prediction, backpropagation, and weight serialization.
 \item \texttt{softmax\_cross\_entropy(logits, labels)} returns mean batch loss and gradient with respect to logits.
 \item \texttt{SGD} updates layer parameters using their stored gradients.
 \item \texttt{build\_shape\_cnn(seed)} builds the fixed untrained Circle/Square/Rectangle architecture.
 \item \texttt{load\_shape\_cnn(path)} loads bundled or caller-provided weights into that architecture.
 \item \texttt{predict\_shape(mask, model)} returns a class label, confidence, and all three probabilities.
 \item \texttt{mycv.streaming.StreamReader} is an optional PyAV-backed source reader; import it explicitly because PyAV is not a core dependency.
\end{itemize}

\section{References}

\begin{enumerate}[label={[\arabic*]}]
\item ITU-R, \emph{Recommendation BT.601-7: Studio Encoding Parameters of
Digital Television}, 2011.
\item C. Harris and M. Stephens, ``A combined corner and edge detector,''
\emph{Proceedings of the Alvey Vision Conference}, pp. 147--151, 1988.
\item R. Hartley and A. Zisserman, \emph{Multiple View Geometry in Computer
Vision}, 2nd ed., Cambridge University Press, 2004.
\item R. C. Gonzalez and R. E. Woods, \emph{Digital Image Processing}, 4th ed.,
Pearson, 2018.
\item G. Bradski, ``The OpenCV Library,'' \emph{Dr. Dobb's Journal of Software
Tools}, 2000. (Background reference; \mycv{} does not depend on OpenCV.)
\item J. P. Lewis, ``Fast normalized cross-correlation,'' \emph{Vision
Interface}, pp. 120--123, 1995.
\item C. E. Shannon, ``Communication in the presence of noise,''
\emph{Proceedings of the IRE}, vol. 37, no. 1, pp. 10--21, 1949.
\item I. Goodfellow, Y. Bengio, and A. Courville, \emph{Deep Learning}, MIT
Press, 2016.
\end{enumerate}

\vfill
\begin{center}\small\color{muted}
Mathematical companion to \mycv{} v4.2.1 \\
Source, tests, model metadata, and release history: \\
\url{https://github.com/charrlie1/mycv2}
\end{center}

\end{document}
